Даны два массива $A$ и $B$ одинакового размера $N$. Сформировать новый массив $C$ того же размера, каждый элемент которого равен максимальному из элементов массивов $A$ и $B$ с тем же индексом.
Дан целочисленный массив $A$ размера $N$. Переписать в новый целочисленный массив $B$ все четные числа из исходного массива (в том же порядке) и вывести размер полученного массива $B$ и его содержимое.
Дан целочисленный массив $A$ размера $N$ ($\leq 15$). Переписать в новый целочисленный массив $B$ все элементы с нечетными порядковыми номерами ($1, 3, \dots$) и вывести размер полученного массива $B$ и его содержимое.
Дан целочисленный массив $A$ размера $N$ ($\leq 15$). Переписать в новый целочисленный массив $B$ все элементы с порядковыми номерами, кратными трем ($3, 6, \dots$), и вывести размер полученного массива $B$ и его содержимое. Условный оператор не использовать.
Дан целочисленный массив $A$ размера $N$. Переписать в новый целочисленный массив $B$ того же размера вначале все элементы исходного массива с четными номерами, а затем --- с нечетными: $A_2, A_4, A_6, \dots, A_1, A_3, A_5, \dots$
Дан массив $A$ размера $N$. Сформировать новый массив $B$ того же размера по следующему правилу: элемент $B_K$ равен сумме элементов массива $A$ с номерами от $1$ до $K$.
Дан массив $A$ размера $N$. Сформировать новый массив $B$ того же размера по следующему правилу: элемент $B_K$ равен среднему арифметическому элементов массива $A$ с номерами от $1$ до $K$.
Дан целочисленный массив размера $N$. Увеличить все четные числа, содержащиеся в массиве, на исходное значение первого четного числа. Если четные числа в массиве отсутствуют, то оставить массив без изменений.
Дан целочисленный массив размера $N$. Увеличить все нечетные числа, содержащиеся в массиве, на исходное значение последнего нечетного числа. Если нечетные числа в массиве отсутствуют, то оставить массив без изменений.
Дан массив размера $N$. Поменять местами его минимальный и максимальный элементы.
Дан массив размера $N$ ($N$ --- четное число). Поменять местами его первый элемент со вторым, третий --- с четвертым и т. д.
Дан массив размера $N$ ($N$ --- четное число). Поменять местами первую и вторую половины массива.
Дан массив размера $N$. Обнулить элементы массива, расположенные между его минимальным и максимальным элементами (не включая минимальный и максимальный элементы).
Дан массив размера $N$. Осуществить сдвиг элементов массива вправо на одну позицию (при этом $A_1$ перейдет в $A_2$, $A_2$ --- в $A_3$, $\dots$, $A_{N - 1}$ --- в $A_N$, a исходное значение последнего элемента будет потеряно). Первый элемент полученного массива положить равным $0$.
Дан массив размера $N$. Осуществить сдвиг элементов массива влево на одну позицию (при этом $A_N$ перейдет в $A_{N - 1}$, $A_{N - 1}$ --- в $A_{N - 2}$, $\dots$, $A_2$ --- в $A_1$, a исходное значение первого элемента будет потеряно). Последний элемент полученного массива положить равным $0$.
Даны два массива $A$ и $B$ одинакового размера $N$. Сформировать новый массив $C$ того же размера, каждый элемент которого равен минимальному из элементов массивов $A$ и $B$ с тем же индексом.
Дан целочисленный массив $A$ размера $N$. Переписать в новый целочисленный массив $B$ все нечетные числа из исходного массива (в том же порядке) и вывести размер полученного массива $B$ и его содержимое.
Дан целочисленный массив $A$ размера $N$ ($\leq 15$). Переписать в новый целочисленный массив $B$ все элементы с четными порядковыми номерами ($2, 4, \dots$) и вывести размер полученного массива $B$ и его содержимое.
Дан целочисленный массив $A$ размера $N$. Сформировать новый целочисленный массив $B$ того же размера, элементы которого вычисляются по правилу: $B_i = A_i \cdot i$ (произведение значения элемента на его порядковый номер).
Дан массив $A$ размера $N$. Сформировать новый массив $B$ того же размера по следующему правилу: элемент $B_K$ равен сумме элементов массива $A$ с номерами от $K$ до $N$.
Дан целочисленный массив размера $N$. Уменьшить все положительные элементы массива на значение первого элемента. Если первый элемент не является положительным, оставить массив без изменений.
Дан массив размера $N$. Заменить все элементы, расположенные до минимального элемента (не включая его), нулями. Если минимальных элементов несколько, обработать первый из них.
Дан массив размера $N$. Поменять местами предпоследний и последний отрицательные элементы массива. Если в массиве меньше двух отрицательных элементов, вывести сообщение об ошибке.
Дан массив размера $N$. Вывести вначале все содержащиеся в массиве четные числа в порядке возрастания их индексов, а затем --- все нечетные числа в порядке убывания их индексов.
Дан целочисленный массив размера $N$. Продублировать (вставить на соседнюю позицию) все элементы, значения которых оканчиваются на заданную цифру $D$. Размер массива при этом увеличивается.
Дан массив размера $N$. Осуществить циклический сдвиг элементов массива вправо на одну позицию (при этом последний элемент перемещается на место первого, а остальные сдвигаются вправо).
Дан массив размера $N$. Сформировать новый массив $B$ из элементов массива $A$, значения которых строго больше среднего арифметического всех элементов $A$.
Дан целочисленный массив размера $N$. Удалить из массива все четные элементы, сжав массив (сдвинув оставшиеся элементы к началу) и изменив его размер.
Дан массив размера $N$ и число $K$ ($1 \leq K \leq N$). Осуществить сдвиг элементов массива влево на $K$ позиций циклически (первые $K$ элементов перемещаются в конец массива, а остальные элементы сдвигаются влево на $K$ позиций).
Дан целочисленный массив размера $N$. Увеличить на 1 все элементы массива, которые больше заданного числа $P$, и уменьшить на 1 все элементы, которые меньше $P$.
